\section{Computing a missing cover margin}
\label{sec:atlas}
A continuous argument often calls a local procedure after choosing a chart containing the point and a neighborhood around it. An effective lower margin allows that chart to be chosen from an approximation. The \repo{Derivations/Lean/IPRADerivations/Foundations/MeshPositivity.lean}{mesh-positivity proof} transfers a supplied positive grid estimate. For finite unions of rational polyhedral charts we can compute the \emph{best} lower slack directly. Its size depends on the chosen inequality presentation.

Let $D=[0,1]^d$. Chart $i$ has a nonempty list of affine forms $\ell_{ij}\in\Q[x_1,\ldots,x_d]$ and accepts a point when all its forms are strictly positive there. Assume at least one chart. Put
\[
 f_i(x)=\min_j\ell_{ij}(x),\qquad f(x)=\max_i f_i(x).
\]
The minimum in the following theorem comes with both a rational value and a rational point of attainment. The internal comparison quantifies over rational codes for $x$.

\begin{thm}[Sharp rational chart margin]\label{th:atlas}
Uniformly in $d$ and in the input tables there are PR outputs $x_0\in D\cap\Q^d$ and $\mu=f(x_0)$ such that every rational $x\in D$ satisfies $f(x)\ge\mu$. The union covers the cube with strict charts exactly when $\mu>0$. In the positive case it also covers regular Cauchy names with a uniform neighborhood radius on which a selected chart works.
\end{thm}
\begin{proof}
Make a list of the $2d$ coordinate face planes and every equality plane $\ell=\ell'$ between input atoms. A constant or zero equation is allowed in the list and can never participate in an independent system of $d$ rows when $d>0$. For every $d$-element sublist, use the determinant to test independence, solve by Cramer's rule, discard solutions outside $D$, and retain the surviving point with least $f$. The cube corners occur and $d=0$ uses the zero-by-zero determinant $1$ and the empty tuple. Every operation is a bounded scan. We show that the list is sufficient without invoking compactness for arbitrary continuous functions.

Take any rational $x\in D$. Sort all atoms by their values at $x$, breaking ties by index. Intersect the cube with the weak inequalities that impose this ordering. On the resulting bounded rational polyhedron $P$, $f$ is one fixed atom $\ell$ because taking the minimum in each group and then the maximum simply selects its position in the order. At $x$, collect the normals of the tight inequalities, including all tight coordinate faces, of rank $r$. If $r=d$, independent tight equations already identify a listed point. If $r<d$, Gaussian elimination gives a rational nonzero vector $v$ annihilated by all these normals. Reverse its sign if needed so that the slope of $\ell$ along it is nonpositive. For each inequality written $a\cdot y\le b$ with $a\cdot v>0$ compute $(b-a\cdot x)/(a\cdot v)$. There is at least one such constraint because the box contains both coordinate faces for each nonzero component of $v$. Each ratio is positive, as tight constraints have zero slope. Move by the smallest of these ratios. The new point is rational, obeys all inequalities, retains every tight equality, and has $f$ no larger than before. The new tight normal with positive pairing against $v$ raises the rank. After at most $d$ moves there is a listed point with $f\le f(x)$. This proves the rational assertion, including lower-dimensional ordering cells and ties.

For each atom let its slope bound be the sum of the absolute values of its linear coefficients, and let $L$ be their maximum. The functions $f_i,f$ are $L$-Lipschitz for the sup norm, since a maximum or minimum of a finite list changes by at most its largest component change. If $\mu\le0$, the rational output $x_0$ fails every strict chart. If $\mu>0$, define $\delta=\mu/(4\max(1,L))$. Given a name in $D$, choose a rational $q\in D$ with error at most $\delta$ by a dyadic precision request and coordinatewise clipping. Choose a maximizing chart for $f(q)$ by rational comparisons. At every point of $D$ within $\delta$ of the named point that same chart has slack at least $\mu-2L\delta\ge\mu/2$. In the arithmetic reading the last statement means the corresponding rational estimates at all requested precisions. Thus one can decide coverage on the supplied finite class and produce a working neighborhood, even when the original specification declared only pointwise positivity.
\end{proof}

\paragraph{Sizes and induction.}
Clear the denominators of \emph{all} affine forms with a common positive multiplier $D_0$. Let $B\ge1$ bound in absolute value every integer coefficient of the resulting arrangement equations, including constants and face equations. The determinants in the enumeration have absolute value at most $d!B^d$ by the permutation formula. Each candidate value of $f$ has a denominator at most $D_0d!B^d$ in lowest terms. In particular, positive $\mu$ has the explicit lower bound $1/(D_0d!B^d)$. The rank computation can scan square minors of size at most $d$. Solving a nonsingular maximal minor with one unmatched coordinate set to one and the other free coordinates zero gives the required null vector. Maximality makes it annihilate the remaining rows as well. To internalize the descent, fix the rational $x$ in the context, run exactly $d$ steps with full-rank states held fixed, and induct on the bounded tests for membership, nonincrease of $\ell$, and $\mathrm{rank}\ge\min(k,d)$. The arithmetic behind Cramer's rule follows by expanding the determinant multilinearly in a column. Every expansion with two equal columns cancels by interchanging them. Rank can equivalently be tracked by elimination with pivot search guarded by a nonzero numerator, an operation on finite tables of rational codes.

\paragraph{Two examples.}
Define four charts of the square by choosing independently between $x<2/3$ or $x>1/3$, and $y<2/3$ or $y>1/3$. With slacks $2/3-x,x-1/3$ and their $y$ analogues, the best common margin is $1/6$, attained for example at $(0,1/2)$. The arrangement includes the tie planes, while testing the original four vertices would yield $2/3$.

For analytic continuation, consider the segment $z(t)=t+i(2t-1)$, $0\le t\le1$, and punctures $0,1$. Each $\|z(t)-a\|_\infty$ is the maximum of four signed affine coordinates. The clearance to the pair is the minimum of two such maxima. Distributing $\min(\max_i u_i,\max_j v_j)=\max_{i,j}\min(u_i,v_j)$ puts it in the required format. Its exact minimum is $1/3$ at $t=1/3,2/3$. This supplies a singularity-clearance input for a path continuation routine. Other inputs, for instance the growth of solutions along the path, require separate arguments. A polygonal path and a finite list of rational punctures work by the same finite distribution and a minimum over segments. The accompanying \path{code/atlas.py} implements the determinant enumeration.
